\section{Item Generation}
\label{sec:itemgen}

% Part 1 begins: the instrument.
The bank is built from two sources: LLM-generated synthetic scenarios and
items extracted from real FDA warning letters. Both share a common schema
and are calibrated identically.

\subsection{Item schema}
Each item is a JSON record with the fields in Table~\ref{tab:schema}.
The \texttt{gold\_standard} is the reference answer used by the judge; for
hard items it must name the controlling regulation at the section level
and, where a Red Herring is present, briefly explain why it is
\emph{not} the violation.

\begin{table}[t]
\centering
\small
\begin{tabular}{@{}ll@{}}
\toprule
field & description \\
\midrule
\texttt{task\_id}      & stable identifier, e.g. \texttt{TASK-21CFR11-A1B2C3D4} \\
\texttt{domain}        & regulatory domain key (see \S\ref{sec:domains}) \\
\texttt{context}       & 4--6 sentence scenario, named fictional company \\
\texttt{question}      & single-answer compliance question \\
\texttt{gold\_standard}& reference answer + section-level citation \\
\texttt{source\_type}  & \texttt{synthetic} or \texttt{real} \\
\texttt{source\_ref}   & URL of the FDA letter (real items only) \\
\bottomrule
\end{tabular}
\caption{Item schema. \texttt{source\_type} tags the provenance track.}
\label{tab:schema}
\end{table}

\subsection{Domains}
\label{sec:domains}
Hard synthetic items cover five regulatory domains---21~CFR~Part~11
\citep{fda21cfr11} (electronic records and signatures), GCP protocol deviations,
promotional-material review, GMP manufacturing deviations, and informed
consent---plus ten \emph{targeted} 21~CFR~Part~11 sub-domains
(hybrid paper/electronic systems, legacy-system validation, audit-trail
completeness, electronic-signature components, and predicate-rule
interplay). Five parallel ``easy'' domains generate deliberately
unambiguous items used to anchor the low-difficulty end of the scale.

\subsection{The Red Herring + subtle violation mandate}
The central design decision is that hard items must be \emph{genuinely}
hard, not merely confusing. The generation prompt enforces two
constraints:
\begin{enumerate}[leftmargin=1.4em,topsep=2pt]
  \item \textbf{Red Herring.} At least one process or element that
  \emph{appears} non-compliant but is in fact fully compliant (e.g.\ a
  legacy system correctly operating read-only, a delayed-but-permissible
  timestamp). It is meant to pull the solver toward flagging a
  non-violation.
  \item \textbf{Subtle violation.} The actual gap must be a procedural
  lapse, a signature-timing issue, or a system-to-system data-integrity
  gap---never a blatant non-compliance. The common-sense answer must be
  \emph{wrong}.
\end{enumerate}
Easy items use a separate template with no Red Herring and a routine,
unambiguous violation. \todo{Add one worked hard-item example (context,
question, gold standard, and the Red Herring) as a boxed figure.}

\subsection{Real-world grounding: FDA warning letters}
A separate importer ingests FDA warning letters, extracts compliance
findings, assigns a domain, and emits items tagged
\texttt{source\_type=real}. Because these encode \emph{actual}
enforcement actions, they serve as an external check on whether the
synthetic difficulty scale is realistic; we return to this in
Section~\ref{sec:results}. \todo{State the final count of real letters
and items in this build.}
